\documentclass[11pt]{article}
\usepackage[margin=1in]{geometry}
\usepackage{amsmath,amssymb}
\newcommand{\finalanswer}[1]{\par\medskip\noindent\fbox{\parbox{0.94\linewidth}{\textbf{Answer:} #1}}}
\title{STAT 412 Homework 6 --- Question 15}
\author{}\date{}
\begin{document}\maketitle
\section*{Problem}
In a survey of $2549$ adults, $1321$ say they have made a New Year's resolution. Construct 90\% and 95\% confidence intervals for the population proportion.
\section*{Work}
\[
\hat p=\frac{1321}{2549}\approx0.5182.
\]
For 90\%, $z=1.645$:
\[
E\approx1.645\sqrt{\frac{0.5182(0.4818)}{2549}}\approx0.0163,
\]
so
\[
(0.5020,0.5345)\approx(0.502,0.535).
\]
For 95\%, $z=1.960$:
\[
E\approx0.0194,\qquad (0.4988,0.5376)\approx(0.499,0.538).
\]

The interpretation is about the population proportion: with the given confidence, the population proportion of adults who say they have made a New Year's resolution is between the endpoints. The 95\% interval is wider because it uses a larger critical value.

\finalanswer{90\% CI $(0.502,0.535)$; 95\% CI $(0.499,0.538)$. Fill-ins: population proportion; between the endpoints. The 95\% interval is wider.}
\end{document}
